\documentclass[runningheads,envcountsame,envcountsect]{llncs}

\let\vec\relax
\usepackage{amsmath,amssymb,mathtools}
\usepackage{booktabs}
\usepackage{graphicx}
\usepackage{microtype}
\usepackage[hidelinks]{hyperref}
\usepackage{cite}
\hypersetup{
  pdftitle={Chosen-Block BAA Codes: Finite-Length Analysis and a Large-Field Rank Obstruction},
  pdfauthor={Stanislav Peceny and Peter Rindal}
}

\newcommand{\F}{\mathbb F}
\newcommand{\wt}{\operatorname{wt}}
\newcommand{\supp}{\operatorname{supp}}
\newcommand{\Acc}{\operatorname{Acc}}
\newcommand{\E}{\mathbb E}
\newcommand{\Prb}{\Pr}
\newcommand{\cC}{\mathcal C}
\newcommand{\cG}{\mathcal G}
\newcommand{\cS}{\mathcal S}

\title{Chosen-Block BAA Codes: Finite-Length Analysis and a Large-Field Rank
Obstruction}
\titlerunning{Chosen-Block BAA Codes}
\author{Stanislav Peceny \and Peter Rindal}
\authorrunning{S. Peceny and P. Rindal}
\institute{}


\begin{document}
\raggedbottom
\maketitle

\begin{abstract}
BAA codes combine a block-diagonal outer code, two interleavers, and two
prefix-sum accumulators.  This paper studies a variant
in which every outer block uses the same certified short code.  Over $\F_2$,
the exact finite-length first moment is obtained by composing the local weight
enumerator with two accumulator kernels.  For extended Golay,
$\operatorname{RM}(2,5)$, and an explicit BCH subcode, numerical optimization
places the relative-distance transition between $95.5\%$ and $99.4\%$ of the
binary Gilbert--Varshamov (GV) bound.  The large-field analysis must instead
group words by support: a single support may carry a large projective family,
which makes the ordinary word enumerator a poor code-generation statistic.
Charging cancellation constraints against shortening dimension gives an
explicit finite union bound.  A backward parity-check matrix also reveals a
coefficient-independent obstruction.  Its generic rank is a layered
path-packing optimum, and the local capacities in that problem can saturate.
For rate-$1/2$ and rate-$1/4$ Reed--Solomon blocks, the numerical structural
transitions are $0.30842$ and $0.60938$, compared with large-field GV values of
$0.49219$ and $0.74358$.  These numbers are diagnostics, not converses for a
typical code.  A Lean 4 formalization checks the finite algebraic and
combinatorial results, the named local codes, and the generic-rank theorem.  It
also exposed three local statement errors, all corrected in the paper; none
affects the finite proof strategy.
\keywords{linear codes \and block--accumulate--accumulate codes \and
minimum distance \and pseudorandom correlation generators \and Lean}
\end{abstract}


\section{Introduction}

Many cryptographic protocols now depend on fast linear codes.  Pseudorandom
correlation generators expand short seeds into OT, VOLE, and multiplication
correlations by multiplying with a structured matrix.  Proximity tests and
code-switching compilers use the same kind of encoder in polynomial commitments
and SNARKs.  At the vector lengths used in these systems, the structure and
constant factors of the encoder matter as much as its asymptotic complexity.

Block--Accumulate codes were introduced by Kolesnikov, Peceny, Rachuri,
Raghuraman, Rindal, and Shah as an implementation-friendly turbo-like family for
PCGs and zero knowledge~\cite{baa26}.  Their main BA-3 construction has the
serial form
\[
    B\Pi_1A\Pi_2A,
\]
where $B$ is a sampled block-diagonal outer map, $\Pi_1,\Pi_2$ are random
interleavers, and $A$ is an accumulator.  That construction comes with detailed
enumeration methods, an asymptotic linear-distance theorem, and efficient CPU
and GPU implementations.  Its field-linear extension uses random nonzero
weighters around the accumulators.  The same paper observes that its local-
distance condition can be enforced by selecting a good subcode and repeating
it in every block.  This paper develops that suggestion at finite length and
asks a second question: over a cryptographic-size field, how close can the
same scalar recursion come to the $q$-ary GV distance?

\paragraph{Results.}
The main results separate finite theorems from asymptotic and numerical claims.

\begin{enumerate}
    \item Starting from the fixed-block specialization suggested in the
    original BAA paper, every map after the repeated local code is made
    explicit as a linear equivalence.  The resulting code has deterministic
    dimension over every finite field, and constant block size gives $O(N)$
    encoding.

    \item Over $\F_2$, the local weight enumerator and binary accumulator
    transition give an exact finite-length first moment.  The best blocks
    considered here reach $96.5\%$ of binary GV at rate $1/2$ and $99.4\%$ at
    rate $1/4$ in the numerical optimization.

    \item Over large fields, exact accumulator traces account for support
    placement, and local shortening dimension accounts for projective freedom.
    The resulting union bound is finite and explicit.

    \item A backward parity-check matrix has generic rank equal to a layered
    path-packing optimum with local block capacities.  The resulting structural
    obstruction holds for every choice of diagonal coefficients.

    \item The finite core is machine checked in Lean 4, including the Golay,
    Reed--Muller, BCH, and Reed--Solomon constructions.  The formalization also
    proves an endpoint-refined container formula.
\end{enumerate}

The two alphabets behave differently.  Binary chosen-block BAA is close to its
GV benchmark.  Large-field BAA has greater absolute symbol distance and smaller
sparse tails, but it remains far from $q$-ary GV.  Since the rank obstruction
does not depend on field values, neither a larger field nor sharper trace
counting can remove it.  Approaching the large-field frontier appears to require
another accumulator layer or a recursion with more state.

\subsection{Technical overview}

The local code enters the two analyses through different summaries.  In the
binary proof, its weight enumerator counts outer words.  Over a larger field,
its shortening dimensions count projective families on a fixed support.  The
direct sum lifts either count to length $N$, and each random interleaver makes
the next accumulator input uniform in its weight or support class.

In the binary case, the support determines every accumulator state.  An input
of weight $w$ produces alternating zero-runs and one-runs, which gives an exact
input--output enumerator.  Composing this enumerator twice yields the expected
number of nonzero outputs below a target weight.  Markov's inequality converts
that expectation into a minimum-distance failure bound.

Over a larger field, the support does not determine the accumulator state.
Several active symbols may cancel, and an $r$-dimensional shortening contains
about $q^{r-1}$ projective candidates.  The finite bound fixes the exact outer
support and both zero-return traces.  It charges only those cancellation
equations that exceed the available projective dimension.

Trace counting still distinguishes many witnesses that lead to the same final
support.  A backward parity-check matrix removes this multiplicity.
After replacing the random diagonal entries by indeterminates, its rank becomes
a two-layer path-packing problem.  Each local MDS block contributes only its
parity-check rank as terminal capacity.  Saturated blocks therefore create
low-weight containers for every specialization of the diagonal entries.

\paragraph{Scope.}
The finite enumerators, probability inequalities, code parameters, and rank
identities are exact.  The entropy limits, sparse-tail estimates, and numerical
transitions are not machine checked.  Neither are the cost model or the claimed
concentration of structural containers.  The paper labels these claims as
derivations, computations, or open problems rather than theorems.

\paragraph{Organization.}
Section~\ref{sec:construction} defines the ensemble.  Sections~\ref{sec:binary}
and~\ref{sec:large-field} give the binary and large-field analyses.
Section~\ref{sec:rank-containers} develops the structural obstruction, and
Section~\ref{sec:numerics} evaluates the concrete local codes.
Section~\ref{sec:implications} discusses design consequences, while
Section~\ref{sec:lean} describes the machine-checked proof.

\subsection{Related work}

\paragraph{Coding theory.}
Turbo codes originated with Berrou, Glavieux, and Thitimajshima~\cite{berrou93}.
Serial concatenation and the uniform-interleaver method were developed and
analyzed by Benedetto, Divsalar, Montorsi, and Pollara~\cite{benedetto98}, while
Divsalar, Jin, and McEliece introduced repeat-accumulate codes and their weight-
enumerator analysis~\cite{divsalar98}.  Kahale and Urbanke studied minimum
distance in random parallel and serial concatenations~\cite{kahale98}.
Bazzi, Mahdian, and Spielman proved that depth-two constant-memory serial
concatenations are asymptotically bad but that repetition followed by two
accumulators can be asymptotically good~\cite{bms09}.  Kliewer, Zigangirov,
Koller, and Costello gave coding theorems for repeat-multiple-accumulate codes,
including punctured and multi-accumulator ensembles approaching binary
GV~\cite{kliewer08}.

Weighted nonbinary RA codes were introduced by Yang~\cite{yang04}; Rosnes and
Graell i Amat derived exact nonbinary accumulator enumerators and asymptotic
distance results for weighted nonbinary repeat-multiple-accumulate
codes~\cite{rosnes11}.  Akhiani and Zhang prove constant relative distance up
to $1/2$ with high probability for RAA codes over sufficiently large finite
fields, improving on the corresponding binary guarantee~\cite{akhiani26}.
Khabbazian obtains sharp minimum-distance lower tails for scaled RAA: when
$r\geq 4$, $0<\delta\leq 1/2$, and $q-1\geq \gamma N$, the failure
probability is $\Theta_{r,\delta,\gamma}(N^{1-r})$~\cite{khabbazian26}.
The projective surplus-cancellation bound used here follows this line of work.
A chosen local block adds a complication: its supports can have many
shortening dimensions, whereas every repetition support has dimension one.

The original BAA work~\cite{baa26} samples systematic block-diagonal local
matrices, computes detailed expected enumerators, proves BA-3 asymptotically
good when each activated block contributes constant distance greater than four,
and gives optimized CPU/GPU implementations.  It also suggests selecting a
good subcode and repeating it across all blocks to enforce the local-distance
condition.  The binary contribution here is an exact finite-length first
moment for each chosen local spectrum, together with certified named blocks.
The large-field contribution is an explicit support-and-shortening analysis
and a generic-rank obstruction for that specialization.  The original paper
also notes an $\F_p$ extension using random coefficients and weighters.  The
monomial ensemble analyzed here accounts explicitly for projective clustering
and structural rank defects.

\paragraph{Cryptographic applications.}
The use of structured LPN and syndrome-decoding assumptions for silent
correlation generation grew from vector-OLE compression and foundational PCG
work~\cite{bcgi18,bcg19pcg}.  Subsequent constructions include Ring-LPN
PCGs~\cite{bcg20ring}, the structured-LDPC-based Silver protocols
~\cite{silver21}, Expand--Accumulate codes~\cite{ea22}, and Expand--Convolute
codes~\cite{ec23}.  Stationary syndrome decoding reduces repeated setup costs
by amortizing the noise-generation protocol~\cite{stationary25}, while recent
work also develops PCGs uniformly across finite fields~\cite{lxyy25}.  BAA was
introduced in this lineage as a particularly hardware-friendly code family
~\cite{baa26}.  Chosen-block BAA preserves the same linear-time architecture
and can therefore be used in the same LPN/SD-style PCG applications.

Code switching and code interleaving make the encoder's rate, distance, and
concrete cost visible in modern proof systems.  Ron-Zewi and Rothblum's constant-overhead IOP and argument framework gives a
theoretical baseline for strictly linear prover computation~\cite{rr25}.
Orion then develops code switching explicitly to obtain a linear-time prover
with polylogarithmic proof size~\cite{orion22}.  Brakedown instantiates a
field-agnostic SNARK with a linear-time encodable code~\cite{brakedown23}.
Block et al.
generalize Expand--Accumulate codes to arbitrary fields and instantiate a
field-agnostic SNARK~\cite{fieldEA24}.  Blaze combines code switching with RAA
codes and supplies improved binary RAA distance bounds~\cite{blaze25}.  Bolt
continues this direction with specially designed sketched codes and lower
commitment constants~\cite{bolt26}.  PIPFRI takes a complementary
``shred-to-shine'' approach, combining linear-time-code proving with a fast
Reed--Solomon FRI layer~\cite{pipfri26}.  Near-GV binary BAA is relevant to
binary code switching in these systems.  Their large-field variants pose the
harder question: can an equally fast scalar encoder approach the $q$-ary
Singleton/GV frontier?

\section{The chosen-block BAA ensemble}
\label{sec:construction}

The ensemble repeats a fixed local code and randomizes only the two recursive
stages.  Its binary and large-field specializations use the same linear maps but
different methods for counting low-weight outputs.

\subsection{Definition and sampling}

Let $q$ be a prime power, let $\F_q$ be the field with $q$ elements, and write
$[N]:=\{1,\ldots,N\}$.  For a finite set $S$, the notation $x\gets S$ means
that $x$ is sampled uniformly from $S$.  Separate sampling statements use
independent randomness unless stated otherwise.  All probabilities and
expectations below are over the displayed ensemble samples; a local subscript
identifies any different probability space.

Let $C_0\subseteq \F_q^m$ be a fixed $[m,k,d]_q$ linear code and let
\[
    B := N/m.
\]
Assume for simplicity that $m\mid N$.  Define the block-diagonal outer code
\[
    C_{\rm out} := C_0^{\oplus B}\subseteq \F_q^N.
\]

Sample independent permutations $\Pi_1,\Pi_2\gets\operatorname{Sym}([N])$.
For $q>2$, let
\[
    D_i=\operatorname{diag}(\alpha_{i,1},\ldots,\alpha_{i,N}),
    \qquad
    \alpha_{i,j}\gets\F_q^\ast,
\]
independently, and set
\[
    M_i := \Pi_i D_i .
\]
For the binary specialization, set $D_i=I$, so $M_i=\Pi_i$.  Permuting before
or after diagonal scaling induces the same nonbinary ensemble.  The formulas
below use the displayed order.

Define the prefix-sum accumulator
\[
    (\Acc x)_j := \sum_{\ell\le j}x_\ell .
\]
The BAA code is
\[
    \boxed{
    \cC_{\rm BAA}
      :=
    \Acc M_2 \Acc M_1(C_{\rm out}).
    }
\]
With row-vector convention, a generator may be written
\[
    G_{\rm BAA}=G_{\rm out}M_1A M_2A,
\]
where $A$ is the lower-triangular all-ones matrix.

\begin{proposition}[Dimension preservation]
For every choice of permutations and nonzero diagonal entries,
\[
    \dim \cC_{\rm BAA} = \frac{k}{m}N
\]
deterministically.
\end{proposition}

\begin{proof}
The accumulator is lower triangular with unit diagonal, and each monomial map
is invertible.  Their composition is a linear equivalence of $\F_q^N$.
Restricting that equivalence to $C_{\rm out}$ preserves dimension, while the
direct sum has dimension $Bk=(k/m)N$.
\end{proof}

\paragraph{Encoding cost.}
For constant $m$, local block encoding costs $O(N)$ field operations.
Each permutation, diagonal scaling, and accumulator also costs $O(N)$.
Hence the complete encoder uses
\[
    \boxed{O(N).}
\]

\subsection{Why two accumulators are necessary}

For comparison, define the one-accumulator code
\[
    \cC_{\rm BA}:=\Acc M_1(C_{\rm out}).
\]

\begin{proposition}[Failure of one accumulation]
Fix a nonzero word in $C_0$ and a constant $0<\delta<1$.  Over the random
monomial map $M_1$, and along multiples of $m$, the one-accumulator ensemble
satisfies
\[
    \Prb[d_{\min}(\cC_{\rm BA})\le \delta N]\longrightarrow 1.
\]
\end{proposition}

\begin{proof}
Fix any nonzero local codeword $c_0\in C_0$ of weight $w$.
Place one copy in each of the $B=N/m$ local blocks.

Let
\[
    H=\{N-L+1,\ldots,N\},
    \qquad L=\lfloor\delta N\rfloor .
\]
For block $i$, let $E_i$ be the event that all $w$ support coordinates of its
copy of $c_0$ land in $H$ under $\Pi_1$.  Then
\[
    \Prb[E_i]
      =
    \frac{\binom{L}{w}}{\binom{N}{w}}
      =
    \delta^w+O(N^{-1}).
\]
On $E_i$, the accumulator input is zero outside $H$, so its output is also
supported in $H$.  Thus $E_i$ produces a codeword of weight at most $L$.

Let $X:=\sum_i\mathbf 1_{E_i}$.  For distinct blocks $i\ne j$, their chosen
supports are disjoint and
\[
    \Prb[E_i\cap E_j]
      =
    \frac{\binom{L}{2w}}{\binom{N}{2w}}
      =
    \Prb[E_i]\Prb[E_j]+O(N^{-1}).
\]
Because $B=\Theta(N)$ and $\Prb[E_i]=\Theta(1)$,
$\E X=\Theta(N)$ and $\operatorname{Var}(X)=O(N)$.  Chebyshev's inequality
therefore gives $\Prb[X=0]=O(N^{-1})$.  Hence
\[
    \Prb[d_{\min}(\cC_{\rm BA})\le \delta N]\to 1.
\]
Thus one accumulator cannot give asymptotically linear distance with a
constant-size direct-sum outer code.  The second randomized accumulator is an
essential part of the construction.
\end{proof}


\section{Binary analysis}
\label{sec:binary}

Set $q=2$.  Then the diagonal randomizers are trivial and the construction is
\[
    \boxed{
    C_0^{\oplus N/m}
      \xrightarrow{\Pi_1}
    \Acc
      \xrightarrow{\Pi_2}
    \Acc .
    }
\]
The local block $C_0\subseteq\F_2^m$ is deterministic.  Its complete weight
spectrum is the only local-code input to the binary analysis.

\subsection{Finite-length first moment}

Write the local homogeneous weight enumerator as
\[
    W_0(z)=\sum_{w=0}^{m}A_w^{(0)}z^w,
    \qquad A_0^{(0)}=1.
\]
For the direct-sum outer code,
\[
    W_{\rm out}(z)=W_0(z)^{N/m},
    \qquad
    A_w^{\rm out}=[z^w]W_0(z)^{N/m}.
    \tag{B1}
\]
For constant $m$, enumerating the $2^k$ local messages gives the complete
spectrum.  A BCH-like block therefore needs neither an asymptotic spectrum
theorem nor a random local-code failure event.

\begin{lemma}[Binary accumulator transition]
For $1\le w\le N$, put
\[
    r_w:=\left\lceil\frac w2\right\rceil,
    \qquad
    s_w:=\left\lfloor\frac w2\right\rfloor.
\]
For $1\le h\le N$, the exact binary accumulator input--output enumerator is
\[
    \boxed{
    E^{\rm bin}_{w,h}
      =
    \binom{h-1}{r_w-1}
    \binom{N-h}{s_w}.
    }
    \tag{B2}
\]
If $w>0$ and $h=0$, then $E^{\rm bin}_{w,0}=0$ because the accumulator is
invertible.  This separate boundary case avoids interpreting
$\binom{h-1}{r_w-1}$ at $h=0$.
\end{lemma}

\begin{proof}
A nonzero input coordinate toggles the accumulator state.  Decomposing the
output into alternating zero-runs and one-runs gives a bijection with two run
compositions.  The positive lengths of the $r_w$ one-runs contribute the first
binomial coefficient.  The $s_w$ required zero gaps contribute the second.
The inverse difference map proves both injectivity and surjectivity of this
correspondence.  Invertibility excludes a zero output for nonzero input.
\end{proof}

Let $\mathcal H_{N,w}:=\{u\in\F_2^N:\wt(u)=w\}$.  Consequently, for
$U\gets\mathcal H_{N,w}$,
\[
    P_N^{\rm bin}(w,h)
      :=
    \Prb_{U\gets\mathcal H_{N,w}}[\wt(\Acc U)=h]
      =
    \frac{E^{\rm bin}_{w,h}}{\binom Nw}.
    \tag{B3}
\]
In contrast with the large-field trace formula below, the binary accumulator
state alternates deterministically between zero and one at the active
coordinates; there is no independently chosen zero-return trace.

Define the second-stage low-weight tail
\[
    Q_N^{\rm bin}(h,L)
      :=
    \sum_{j=1}^{L}P_N^{\rm bin}(h,j).
\]
\begin{theorem}[Exact binary first moment]
Let $1\le L\le N$.  Over independent
$\Pi_1,\Pi_2\gets\operatorname{Sym}([N])$, the expected number of nonzero BAA
codewords of weight at most $L$ is
\[
    \boxed{
    \mu_{\rm bin}(N,L;C_0)
      =
    \sum_{w=1}^{N}A_w^{\rm out}
    \sum_{h=1}^{N}
      P_N^{\rm bin}(w,h)
      Q_N^{\rm bin}(h,L).
    }
    \tag{B4}
\]
Both interleavers are uniform, so conditioned on its weight, the word entering
each accumulator is uniform on the corresponding Hamming slice.  Since the
permutations and accumulators are invertible,
\[
    \boxed{
    \Prb[d_{\min}(\cC_{\rm BAA})\le L]
      \le
    \mu_{\rm bin}(N,L;C_0).
    }
    \tag{B5}
\]
Unlike the large-field case, no projective-support correction is needed:
$\F_2^\ast$ contains only one scalar.
\end{theorem}

\begin{proof}
Fix an outer word of weight $w$.  A uniform permutation sends it uniformly to
the weight-$w$ Hamming slice, so (B3) gives the first output weight.  Conditional
on an intermediate weight $h$, the independent second permutation is uniform
on the weight-$h$ slice, and $Q_N^{\rm bin}(h,L)$ gives the final tail.
Summing over the $A_w^{\rm out}$ outer words proves (B4) by linearity of
expectation.  If the minimum distance is at most $L$, at least one nonzero bad
word exists.  Markov's inequality therefore gives (B5).
\end{proof}

\subsection{Asymptotic consequences}

Fix $w$ and a target $0<\delta<1/2$.  The exact transition~(B3) gives, for
fixed $h$,
\[
    P_N^{\rm bin}(w,h)
      =
    \Theta_{w,h}\!\left(N^{-\lceil w/2\rceil}\right).
\]
The second accumulator tail from fixed $h$ converges to a nonzero beta-tail
constant.  Splitting the intermediate-weight sum into bounded and growing
ranges therefore gives
\[
    \boxed{
    \sum_h P_N^{\rm bin}(w,h)Q_N^{\rm bin}(h,\lfloor\delta N\rfloor)
      =
    \Theta_{w,\delta}\!\left(N^{-\lceil w/2\rceil}\right).
    }
    \tag{B6}
\]
Hence the one-active-block contribution is
\[
    \frac Nm
    \sum_{w=1}^{m}A_w^{(0)}
    \Theta_{w,\delta}\!\left(N^{-\lceil w/2\rceil}\right).
    \tag{B7}
\]
If the local distance is $d$, this is at most
\[
    O\!\left(N^{1-\lceil d/2\rceil}\right),
\]
with the exact leading term determined by all local weights attaining the
smallest value of $\lceil w/2\rceil$.

Using only local distance, a fixed family touching $t$ blocks has the coarse
scale
\[
    O\!\left(N^{t-\lceil td/2\rceil}\right).
    \tag{B8}
\]
Thus a constant local distance of at least three is consistent with a vanishing
fixed-$t$ sparse tail.  The earlier binary BA proof assumes $d>4$ because it
uses coarser inequalities.  For a fixed local block, its exact spectrum and the
two-accumulator kernel~(B4) give the sharper bound.

\paragraph{Linear weights.}

The direct-sum spectrum has the coefficient upper exponent
\[
    \boxed{
    \mathcal A_{C_0}(\alpha)
      :=
    \inf_{z>0}
    \left\{
      \frac1m\log W_0(z)-\alpha\log z
    \right\}.
    }
    \tag{B9}
\]
Indeed,
\[
    A_{\lfloor\alpha N\rfloor}^{\rm out}
      \le
    \exp\!\left(N\mathcal A_{C_0}(\alpha)+O(\log N)\right).
\]
For linear input and output weights $uN$ and $vN$, the binary accumulator
transition exponent is
\[
    \boxed{
    \Gamma_{\rm bin}(u,v)
      :=
    vH\!\left(\frac{u}{2v}\right)
    +(1-v)H\!\left(\frac{u}{2(1-v)}\right)
    -H(u),
    }
    \tag{B10}
\]
with feasibility $u/2\le v\le1-u/2$.  For $\delta<1/2$, this exponent is
increasing in the final output fraction up to $1/2$, so the low-weight tail is
maximized at the boundary $v=\delta$.

The binary BAA first-moment exponent is therefore
\[
    \boxed{
    \Phi^{\rm bin}_{C_0,\delta}(\alpha,\eta)
      =
    \mathcal A_{C_0}(\alpha)
    +\Gamma_{\rm bin}(\alpha,\eta)
    +\Gamma_{\rm bin}(\eta,\delta).
    }
    \tag{B11}
\]
A sufficient condition for the linear-weight contribution to vanish is
\[
    \sup_{\alpha,\eta}\Phi^{\rm bin}_{C_0,\delta}(\alpha,\eta)<0
    \tag{B12}
\]
over the feasible domain.  For a fixed block $C_0$, (B4) gives the exact
finite-length first moment, (B7) controls fixed-block supports, and (B12)
controls linear weights.

\section{Large-field support analysis}
\label{sec:large-field}

Over a large field, one support can carry many scalar assignments.  Counting
them as separate bad words makes the first moment sensitive to large projective
clusters.  The bound below fixes the support first and charges only the
cancellation constraints beyond its projective dimension.

\subsection{Accumulator traces and projective cancellation}

The use of random nonzero weighters around nonbinary accumulators follows the
weighted nonbinary RA/RMA literature~\cite{yang04,rosnes11}.  The distinction
between counting field-valued codewords and bounding the probability that a
sampled code admits a bad projective word is central to recent large-field RAA
analyses~\cite{akhiani26,khabbazian26}.  In chosen-block BAA, the shortening
lattice of the local code replaces the one-dimensional supports of the
repetition code.


Consider an accumulator input with exactly $s$ nonzero coordinates at
\[
    p_1<\cdots<p_s,
    \qquad p_{s+1}:=N+1.
\]
Let
\[
    z_i:=\sum_{\ell=1}^i u_{p_\ell}
\]
be the accumulator state immediately after the $i$th active coordinate.
Define the zero-return trace
\[
    Z:=\{i:z_i=0\},
    \qquad a:=|Z|.
\]
Because every active input symbol is nonzero, $1\notin Z$ and $Z$ cannot contain
consecutive indices.

\begin{lemma}[Nonbinary accumulator placements]
\label{lem:nonbinary-placement}
Suppose the active input coordinates form a uniform $s$-subset of $[N]$.  The
support-placement probability for output weight $h$, summed over all
zero-return traces of cardinality $a$, is
\[
    \boxed{
    \cG_N(s,h,a)
      :=
    \frac{
      \binom{s-a}{a}
      \binom{N-h}{a}
      \binom{h-1}{s-a-1}
    }{
      \binom{N}{s}
    }.
    }
    \tag{1}
\]
for feasible $s,h,a$; the value is zero otherwise.
\end{lemma}

\begin{proof}
The three factors count, respectively:
\begin{itemize}
    \item the nonadjacent zero-return trace;
    \item the zero-state gaps;
    \item the positive nonzero-state gaps.
\end{itemize}
Dividing their product by the $\binom Ns$ possible input supports gives (1).
The gap construction is reversible, so every placement is counted once.
\end{proof}

Fix an exact outer support $S\subseteq[N]$.  Write
\[
    s:=|S|,
    \qquad
    C_{\rm out}[S]
      :=
    \{c\in C_{\rm out}:\supp(c)\subseteq S\},
    \qquad
    r:=\dim C_{\rm out}[S].
\]
The shortening contains exactly
\[
    P_r=\frac{q^r-1}{q-1}
\]
projective lines, of which at most $P_r$ can have exact support $S$.

For one fixed projective word, a prescribed zero return in one accumulator
forces one fresh diagonal multiplier to one value in $\F_q^\ast$.
Hence a prescribed total of $c$ zero returns across the two stages succeeds
with probability at most $(q-1)^{-c}$ for that fixed projective word.

\begin{lemma}[Projective cancellation bound]
\label{lem:projective-cancellation}
Let $q>2$, let an outer shortening have dimension $r\ge1$, and prescribe $c$
zero returns across the two accumulators.  Over the independent nonzero
diagonal entries, the probability that some projective word in the shortening
realizes all $c$ returns is at most
\[
    \boxed{
    \chi_q(r,c)
      :=
    \min\left\{
        1,\;
        \frac{q^r-1}{q-1}(q-1)^{-c}
    \right\}.
    }
    \tag{2}
\]
\end{lemma}

\begin{proof}
Normalize each nonzero word by its first nonzero coordinate.  An
$r$-dimensional shortening has $(q^r-1)/(q-1)$ normalized representatives.
For a fixed representative, each prescribed return fixes a fresh independent
diagonal entry and therefore costs a factor $(q-1)^{-1}$.  The union bound over
representatives gives the second argument of the minimum in (2); probability
is at most one.
\end{proof}

For large $q$,
\[
    \chi_q(r,c)
      \approx
    (q-1)^{-(c-(r-1))_+}.
\]
Thus the first $r-1$ cancellation constraints can be paid for by projective
degrees of freedom; each surplus constraint costs approximately $1/q$.

\subsection{Finite-length master bound}

\begin{theorem}[Finite large-field master bound]
Let $q>2$ and $1\le L\le N$.  Over independent uniform permutations and
independent uniform nonzero diagonal entries, the chosen-block BAA ensemble
satisfies
\[
\boxed{
\begin{aligned}
\Prb[d_{\min}(\cC_{\rm BAA})\le L]
\le
\sum_{\substack{\varnothing\ne S\subseteq[N]}}
\sum_{a=0}^{|S|-1}
\sum_{h=1}^{N}
\sum_{b=0}^{h-1}
\sum_{j=1}^{L}
&\cG_N(|S|,h,a)
\\[-1mm]
&\cdot \cG_N(h,j,b)
\cdot \chi_q(r_S,a+b),
\end{aligned}
}
\tag{3}
\]
where $r_S:=\dim C_{\rm out}[S]$.
\end{theorem}

\begin{proof}
Partition a nonzero outer word by its exact support $S$, and partition each
accumulator trajectory by $(a,h,b,j)$.  Lemma~\ref{lem:nonbinary-placement}
supplies the two independent support-placement factors.  For each fixed
support, Lemma~\ref{lem:projective-cancellation} bounds the union over its
normalized projective words by $\chi_q(r_S,a+b)$.  A final union bound over the
displayed finite index set proves (3).
\end{proof}

Although $S$ indexes an exact-support event, $r_S$ is the dimension of the
larger shortening whose words have support contained in $S$.  It therefore
upper-bounds the number of projective candidates with exact support $S$.
Equation~(3) groups each shortened projective space before applying the union
bound, so every codeword retains its unique exact support.

\subsection{Shortenings and asymptotic consequences}

For a local coordinate set $T\subseteq[m]$, define
\[
    C_0[T]
      :=
    \{c\in C_0:\supp(c)\subseteq T\},
    \qquad
    r_0(T):=\dim C_0[T].
\]

Call $T$ \emph{admissible} when $C_0[T]$ is nontrivial.

\begin{proposition}[Shortening profile]
Suppose that $C_0$ has minimum distance $d$.  Every admissible set $T$ of size
$s_T$ satisfies
\[
    s_T-r_0(T)\ge d-1.
\]
Consequently, let a global support have size $s$, shortening dimension $r$,
and nonempty intersection with $t$ blocks.  If each of those $t$ local
intersections is admissible, then
\[
    s-r\ge t(d-1),
\]
and its support excess $E:=s-r+1$ satisfies
\[
    \boxed{
    E\ge t(d-1)+1.
    }
    \tag{4}
\]
If $C_0$ is an $[m,k,d]$ MDS code, then every $T\subseteq[m]$ satisfies
\[
    \boxed{
    r_0(T)=(|T|-d+1)_+.
    }
    \tag{5}
\]
\end{proposition}

\begin{proof}
The shortening $C_0[T]$ has length $s_T$, dimension $r_0(T)$, and minimum
distance at least $d$.  Singleton gives the local inequality.  Shortening of a
direct sum decomposes blockwise, so dimensions and support sizes add; summing
over the $t$ admissible intersections proves (4).

For an MDS code, restriction to any $k$ coordinates is an isomorphism.  The
kernel of restriction to the complement of $T$ therefore has dimension
$\max\{0,k-(m-|T|)\}=(|T|-d+1)_+$, which proves (5).
\end{proof}

The admissibility condition in (4) is necessary: an arbitrary small set can
support only the zero shortening and need not satisfy the local Singleton
inequality.  Exact supports of outer codewords are admissible in every block
they meet.

This pointwise optimal shortening profile makes a small Reed--Solomon block a
natural local code over a large field.

\paragraph{Sparse supports.}

Suppose an admissible global support touches $t$ local blocks.
In the no-surplus regime,
\[
    a+b\le r-1.
\]
The first accumulator output has at least $s-a$ nonzero-state gaps, and after
the second stage any final output of weight $j$ satisfies
\[
    j\ge s-a-b.
\]
Using the full projective cancellation budget gives
\[
    s-a-b\ge s-r+1\ge t(d-1)+1.
\]

A direct fixed-$t$ asymptotic analysis of the exact kernel gives the scale
\[
    \boxed{
    P_t
      =
    O\!\left(
      N^{-t(d-2)-1}
    \right)
    }
    \tag{6}
\]
when $q$ is sufficiently large that surplus cancellations are negligible.

In particular,
\[
    P_1=O(N^{1-d}).
\]
The canonical one-block event gives a matching lower bound of the same
polynomial order, so the natural sparse tail is
\[
    \boxed{
    \Prb[\text{sparse bad event}]
      =
    \Theta(N^{1-d}).
    }
    \tag{7}
\]

For example:
\[
\begin{array}{c|c}
d & \text{sparse tail scale}\\
\hline
2 & N^{-1}\\
3 & N^{-2}\\
5 & N^{-4}\\
7 & N^{-6}\\
9 & N^{-8}
\end{array}
\]

For $q-1\gtrsim N$, a surplus cancellation can recover at most one geometric
power of $N$ while costing $1/(q-1)$, so surplus layers remain controlled.
For $q\approx2^{128}$ and the lengths considered here, the resulting surplus
terms are negligible on the reported scale.

\paragraph{Linear weights.}

Let $C_0$ be an $[m,k,d]$ MDS code, with
\[
    d=m-k+1,
    \qquad R=k/m.
\]
Its local support--dimension polynomial is
\[
    \boxed{
    F_{m,k}(z,u)
      =
    1+
    \sum_{\ell=d}^{m}
    \binom{m}{\ell}z^\ell u^{\ell-d+1}.
    }
    \tag{L1}
\]
Let $p_0,p_d,\ldots,p_m$ be the fractions of local blocks in each exact
support class, and put
\[
    \lambda:=\sum_{\ell=d}^{m}p_\ell.
\]
Define the normalized global support and shortening dimension by
\[
    \alpha
      :=
    \frac1m\sum_{\ell=d}^{m}\ell p_\ell,
    \qquad
    \beta
      :=
    \frac1m\sum_{\ell=d}^{m}(\ell-d+1)p_\ell.
\]
Since $d-1=m-k=m(1-R)$,
\[
    \boxed{
    \beta=\alpha-(1-R)\lambda.
    }
    \tag{L2}
\]
The support-pattern entropy per output coordinate is
\[
    \boxed{
    \cS_{m,k}(p)
      :=
    \frac1m
    \left[
      H(p_0,p_d,\ldots,p_m)
      +
      \sum_{\ell=d}^{m}p_\ell\log\binom{m}{\ell}
    \right].
    }
    \tag{L3}
\]
Here $H$ uses natural logarithms unless a quantity is explicitly divided by
$\log 2$.

Normalize the first- and second-stage zero-return counts and intermediate
support by
\[
    x:=a/N,
    \qquad
    y:=b/N,
    \qquad
    \eta:=h/N.
\]
The first accumulator contributes
\[
\boxed{
\begin{aligned}
\Gamma_1(\alpha,\eta,x)
={}&
(\alpha-x)
H\!\left(\frac{x}{\alpha-x}\right)
+
(1-\eta)
H\!\left(\frac{x}{1-\eta}\right)
\\
&+
\eta
H\!\left(\frac{\alpha-x}{\eta}\right)
-H(\alpha),
\end{aligned}
}
\tag{L4}
\]
and, at final relative weight $\zeta$, the second contributes
\[
\boxed{
\begin{aligned}
\Gamma_2(\eta,\zeta,y)
={}&
(\eta-y)
H\!\left(\frac{y}{\eta-y}\right)
+
(1-\zeta)
H\!\left(\frac{y}{1-\zeta}\right)
\\
&+
\zeta
H\!\left(\frac{\eta-y}{\zeta}\right)
-H(\eta).
\end{aligned}
}
\tag{L5}
\]
A direct derivative calculation gives the unconstrained maximizing final weight
\[
    \boxed{
    \zeta_\star=1-\frac{y}{\eta}.
    }
    \tag{L6}
\]
Therefore define
\[
    \widehat\Gamma_2(\eta,\delta,y)
      :=
    \Gamma_2\!\left(
      \eta,
      \min\left\{\delta,1-\frac{y}{\eta}\right\},
      y
    \right).
\]
This distinction matters for target distances exceeding $1/2$.

In the large-field no-surplus region,
\[
    x+y\le\beta,
\]
and the complete support/trace first-moment exponent is
\[
    \boxed{
    \Phi_{m,k,\delta}
      =
    \cS_{m,k}(p)
      +
    \Gamma_1(\alpha,\eta,x)
      +
    \widehat\Gamma_2(\eta,\delta,y).
    }
    \tag{L7}
\]
For finite $q$, the corresponding exponent includes
\[
    \min\left\{
      0,
      \beta\log q-(x+y)\log(q-1)
    \right\}.
\]
The numerical transition for this support/trace bound is the largest $\delta$
for which the supremum of~(L7) over the feasible domain is negative.

\section{Generic rank and the structural obstruction}
\label{sec:rank-containers}

The master bound in Section~\ref{sec:large-field} distinguishes zero-return
traces even when they produce the same final support.  Its first moment is
therefore not a code-level converse.  A backward parity-check matrix groups all
traces and projective assignments that lead to one support.  The generic rank
of that matrix has an exact path-packing characterization and yields a
coefficient-independent obstruction.

\subsection{Backward parity-check formulation}

Switch temporarily to column-vector convention.  Let $H$ be a parity-check
matrix for $C_{\rm out}$ and let
\[
    \Delta:=A^{-1}
\]
be the first-difference operator.  After renaming the inverses of the two random
monomial maps, a parity-check matrix for BAA is
\[
    \boxed{
    Q=H D_1P_1\Delta D_2P_2\Delta .
    }
    \tag{G1}
\]
For a final coordinate set $T\subseteq[N]$, a nonzero BAA codeword supported in
$T$ exists exactly when
\[
    \boxed{\operatorname{rank}Q_T<|T|.}
    \tag{G2}
\]
This event groups every outer support, projective assignment, and pair of
accumulator traces that leads to the same final support.

Replace the entries of $D_1$ and $D_2$ by two disjoint families of
indeterminates, and let $K$ be the fraction field of the resulting polynomial
ring.  All ranks in the remainder of this subsection are over $K$.  For an
$[m,k]$ MDS block, define the local capacity
\[
    h:=m-k
\]
for the number of paths that may terminate in that block.

The layered graph associated with (G1) has three copies $V_2,V_1,V_0$ of
$[N]$.  Add an edge from $j\in V_2$ to $i\in V_1$ when
$(P_2\Delta)_{i,j}\ne0$, and add an edge from $j\in V_1$ to $i\in V_0$ when
$(P_1\Delta)_{i,j}\ne0$.  Join each vertex of $V_0$ to the local block that
contains its coordinate.  Coordinate vertices have unit capacity, and every
block terminal has capacity $h$.  For $T\subseteq V_2$, let $\nu(T)$ be the
maximum number of capacity-respecting paths from distinct vertices of $T$ to
the block terminals.

\begin{theorem}[Generic rank as path packing]
Let $C_0$ be an $[m,k]$ MDS code, and let $H$ be a block-diagonal parity-check
matrix for $C_{\rm out}$.  For every final support $T\subseteq[N]$,
\[
    \operatorname{rank}_K Q_T=\nu(T).
\]
Equivalently, the generic rank equals the minimum cut capacity in the layered
graph above.
\end{theorem}

\begin{proof}
Apply rectangular Cauchy--Binet at both diagonal layers.  Each pair of selected
coordinate sets contributes a distinct diagonal monomial, so nonzero terms
cannot cancel across selections.  A minor of $\Delta$ is nonzero exactly when
its selected row and column sets admit the unique vertex-disjoint linkage in
which each step stays in place or moves to its successor.  Finally, an MDS
parity-check matrix has selected-column rank $\min\{s,h\}$ on a block with $s$
selected columns.  A full-rank minor therefore exists exactly when the layered
network admits a capacity-respecting path packing of the same size.
\end{proof}

This theorem is a trace-free analogue of Khabbazian's shared-rank method for
large-field RAA~\cite{khabbazian26}.

\subsection{A coefficient-independent obstruction}

For $S\subseteq[N]$, define
\[
    S+1:=\{i+1:i\in S,\ i<N\}
\]
and the one-step path neighborhood
\[
    \partial S:=S\cup(S+1),
\]
so the right endpoint is handled explicitly.  Define the two-step backward
neighborhood
\[
    U(T):=P_1\partial\!\left(P_2\partial T\right).
    \tag{G3}
\]
For an outer-coordinate set $U$, define its MDS block rank
\[
    \kappa(U):=
    \sum_{\text{blocks }B}\min\{|U\cap B|,h\}.
    \tag{G4}
\]
Every backward path from $T$ terminates in $U(T)$, hence
\[
    \operatorname{rank}Q_T\le \kappa(U(T)).
\]
Consequently, the rank criterion gives the following corollary.

\begin{proposition}[Structural obstruction]
Let $C_0$ be an $[m,k]$ MDS code and put $h=m-k$.  For every
$T\subseteq[N]$,
\[
    \boxed{
    \kappa(U(T))<|T|
    \quad\Longrightarrow\quad
    T\text{ supports a nonzero BAA codeword}
    }
    \tag{G5}
\]
for every choice of diagonal coefficients.  The obstruction is purely a defect
of the two permutations, the two difference layers, and the local block-rank
caps.  Increasing $q$ cannot remove it.
\end{proposition}

\begin{proof}
The selected columns of the transported parity-check matrix lie in the direct
sum of the blockwise column spaces indexed by $U(T)$.  That space has dimension
at most $\kappa(U(T))$.  If this value is smaller than $|T|$, then $Q_T$ has a
nonzero kernel vector for every specialization of the diagonal entries.  By
(G2), the resulting codeword is supported in $T$.
\end{proof}

The mechanism is block saturation.  Once $U(T)$ occupies at least $h$ positions
inside a block, additional positions in that block add no parity-check rank.
Two backward neighborhood expansions can therefore contain many coordinates
but concentrate them in a smaller number of almost-full blocks.

\subsection{Counting structural containers}

For fixed permutations, define
\[
    Z_t
      :=
    \bigl|\{T\subseteq[N]:|T|=t,\ \kappa(U(T))<t\}\bigr|.
\]
By (G5), every set counted by $Z_t$ supports a nonzero BAA codeword for every
choice of diagonal coefficients.

Endpoint membership changes the size of a difference neighborhood.  For
$e\in\{0,1\}$, let $R_N^e(t,r)$ count the $t$-subsets with $r$ runs and with
$e=1$ exactly when $N$ belongs to the set.  Then
\[
\begin{aligned}
    R_N^0(t,r)
      &:=\binom{t-1}{r-1}\binom{N-t}{r},\\
    R_N^1(t,r)
      &:=\binom{t-1}{r-1}\binom{N-t}{r-1},\\
    R_N(t,r)
      &:=R_N^0(t,r)+R_N^1(t,r)
       =\binom{t-1}{r-1}\binom{N-t+1}{r}.
\end{aligned}
\tag{G6}
\]
For a set of type $(t,r,e)$, the neighborhood has exact size
\[
    |\partial T|=t+r-e.
\]

For the local block-rank count define
\[
    B_{m,h}(x,y)
      :=
    \sum_{\ell=0}^{m}\binom m\ell x^\ell y^{\min(\ell,h)},
    \tag{G7}
\]
and
\[
    A_{N,m,h}(u,c)
      :=
    \sum_{z\le c}[x^uy^z]B_{m,h}(x,y)^{N/m}.
    \tag{G8}
\]

\begin{theorem}[Exact structural-container expectation]
Assume $m\mid N$, fix $1\le t\le N$, and sample independent
$P_1,P_2\gets\operatorname{Sym}([N])$.  For
$r_1,r_2\ge1$ and $e_1,e_2\in\{0,1\}$, set
\[
    u_1:=t+r_1-e_1,
    \qquad
    u_2:=u_1+r_2-e_2.
\]
Let $\mathcal I_t$ contain the tuples $(r_1,r_2,e_1,e_2)$ that satisfy
$r_1\le t$, $r_2\le u_1$, and $u_1,u_2\le N$.
Then the expected number of direct structural containers $T$ of size $t$ is
\[
\boxed{
\begin{aligned}
\E Z_t
=
\sum_{(r_1,r_2,e_1,e_2)\in\mathcal I_t}
&R_N^{e_1}(t,r_1)
\frac{R_N^{e_2}(u_1,r_2)}{\binom N{u_1}}
\\[-1mm]
&\cdot
\frac{A_{N,m,h}(u_2,t-1)}
     {\binom N{u_2}}.
\end{aligned}
}
\tag{G9}
\]
\end{theorem}

\begin{proof}
Partition the $t$-subsets by their first run count and endpoint flag.  Equation
(G6) gives $R_N^{e_1}(t,r_1)$ choices, and the first neighborhood has size
$u_1$.  The permutation $P_2$ maps this neighborhood to a uniform
$u_1$-subset.  A second application of (G6) gives the ratio
$R_N^{e_2}(u_1,r_2)/\binom N{u_1}$ and a neighborhood of size $u_2$.  The
permutation $P_1$ maps that neighborhood to a uniform $u_2$-subset.  Equation
(G8), divided by $\binom N{u_2}$, is exactly the probability that its local
block rank is at most $t-1$.  Multiplying the factors and summing the disjoint
cases proves (G9).
\end{proof}

Impossible parameter choices contribute zero under the usual binomial
convention.  Equation~(G9) counts each final support once and contains no field
values or zero-return choices.  Omitting $e_1,e_2$ changes neighborhood sizes
by at most two and gives the same entropy exponent, but not the same finite
expectation.

\section{Numerical evaluation}
\label{sec:numerics}

The calculations below evaluate the binary first moment, the large-field trace
bound, and the structural-container count.  All transition points come from
floating-point optimization; none is an interval-certified theorem.

\subsection{Binary instantiations}

The binary Gilbert--Varshamov benchmarks solve
\[
    H_2(\delta_{\rm GV})=1-R.
\]
Numerically,
\[
    \delta_{\rm GV}(R=1/2)\approx0.1100278644,
    \qquad
    \delta_{\rm GV}(R=1/4)\approx0.2145017449.
\]

The short-code parameters and enumerators are standard; see MacWilliams and
Sloane~\cite{macwilliams}.  For each local spectrum, the calculation maximizes
(B11) over its feasible domain and finds the value of $\delta$ where the
nontrivial interior maximum crosses zero.  Equations (B6)--(B8) handle sparse
weights separately.

\paragraph{Rate $1/2$: extended Golay.}
The extended binary Golay code has parameters $[24,12,8]$ and spectrum
\[
    W_{\rm Golay}(z)
      =
    1+759z^8+2576z^{12}+759z^{16}+z^{24}.
\]
The numerical BAA transition is
\[
    \boxed{\delta^\star_{\rm Golay}\approx0.1051193371.}
\]
The contribution from one active block has order $N^{-3}$.

\paragraph{Rate $1/2$: $\operatorname{RM}(2,5)$.}
The second-order Reed--Muller code has parameters $[32,16,8]$ and spectrum
\[
\begin{aligned}
    W_{\operatorname{RM}(2,5)}(z)
      ={}&
    1+620z^8+13888z^{12}+36518z^{16}
    \\
    &+13888z^{20}+620z^{24}+z^{32}.
\end{aligned}
\]
The numerical BAA transition is
\[
    \boxed{\delta^\star_{\operatorname{RM}(2,5)}\approx0.1061649861.}
\]
Among the rate-$1/2$ blocks in this evaluation, $\operatorname{RM}(2,5)$ gives
the largest transition: about $96.5\%$ of the binary GV relative distance.  The
one-block contribution again has order $N^{-3}$.

\paragraph{Rate $1/4$: a concrete $[32,8,12]$ BCH subcode.}
Start from the extended primitive BCH code $[32,11,12]$ and select the explicit
$8$-dimensional subcode given in Appendix~\ref{app:binary-generator}.  Direct
enumeration of its $256$ words gives
\[
    \boxed{
    W_{32,8,12}(z)
      =
    1+48z^{12}+158z^{16}+48z^{20}+z^{32}.
    }
\]
The corresponding numerical BAA transition is
\[
    \boxed{\delta^\star_{32,8,12}\approx0.2132933411.}
\]
This is about $99.4\%$ of the binary rate-$1/4$ GV distance.  The contribution
from one active block has order $N^{-5}$.

Fixing the short code makes its complete spectrum available to both the
optimization and an exhaustive check.  The global encoding cost remains
linear in $N$.

\subsection{Large-field benchmarks and method}

For a $q$-ary linear code of rate $R$, the Gilbert--Varshamov relative distance
solves $H_q(\delta_{\rm GV})=1-R$, where
\[
    H_q(\delta)
      :=
    \delta\log_q(q-1)
    -\delta\log_q\delta
    -(1-\delta)\log_q(1-\delta).
\]
For $q=2^{128}$, the rate-$1/2$ and rate-$1/4$ benchmarks are
\[
    \delta_{\rm GV}(1/2)\approx0.492189,
    \qquad
    \delta_{\rm GV}(1/4)\approx0.743584.
    \tag{R2}
\]
The corresponding Singleton ceilings are $0.5$ and $0.75$.  For each local
block below, the calculation maximizes (L7) over its feasible domain and finds
the distance at which the maximum crosses zero.

\subsection{Large-field instantiations}

\paragraph{Rate $1/2$.}

Take
\[
    C_0=\operatorname{RS}[16,8,9].
\]
Then
\[
    F_{16,8}(z,u)
      =
    1+
    \sum_{\ell=9}^{16}
    \binom{16}{\ell}z^\ell u^{\ell-8},
\]
and
\[
    \beta=\alpha-\frac12\lambda.
\]
Numerical optimization of the support/trace exponent gives the transition
\[
    \boxed{
    \delta^\star_{16,8}\approx0.3052795.
    }
    \tag{R1}
\]
At the transition, a representative saddle is approximately
\[
\begin{aligned}
    \lambda &\approx0.5000, &
    \alpha &\approx0.4477, &
    \beta &\approx0.1977,\\
    x &\approx0.1076, &
    y &\approx0.0901, &
    \eta &\approx0.3670.
\end{aligned}
\]
At this saddle, the available projective dimension is almost exactly exhausted
by the zero-return constraints:
\[
    \boxed{x+y\approx\beta.}
\]
At $\delta=0.30$, the optimized exponent is approximately
\[
    \Phi_{16,8,0.30}^{\max}/\log 2\approx-0.016699
\]
bits per output coordinate.

The local distance is $d=9$, so the sparse one-block scale is $N^{-8}$.  At
output length $N=2^{20}$ and $\delta=0.30$, the marked one-block term is
approximately $2^{-129.57}$.  The sparse term and the negative linear-weight
exponent are consistent with $d_{\min}>0.30N$, but do not prove that inequality
at finite length.

\paragraph{Rate $1/4$.}

Take
\[
    C_0=\operatorname{RS}[16,4,13].
\]
The local support--dimension polynomial is
\[
    \boxed{
    F_{16,4}(z,u)
      =
    1+560z^{13}u+120z^{14}u^2+16z^{15}u^3+z^{16}u^4.
    }
    \tag{R3}
\]
With
\[
    \lambda=p_{13}+p_{14}+p_{15}+p_{16},
\]
equation~(L2) becomes
\[
    \boxed{
    \beta=\alpha-\frac34\lambda.
    }
    \tag{R4}
\]
The larger redundancy sharply reduces the projective cancellation budget.

Numerical optimization of~(L7) gives the estimated transition
\[
    \boxed{
    \delta^\star_{16,4}\approx0.60441329.
    }
    \tag{R5}
\]
At the transition, a representative saddle is approximately
\[
\begin{aligned}
    \lambda &\approx0.6701545, &
    \alpha &\approx0.6601815, &
    \beta &\approx0.1575656,\\
    x &\approx0.0648292, &
    y &\approx0.0927364, &
    \eta &\approx0.6655770.
\end{aligned}
\]
Here too, the zero-return constraints nearly exhaust the projective dimension:
\[
    \boxed{x+y\approx\beta.}
\]
At this saddle,
\[
    1-\frac{y}{\eta}\approx0.8607>\delta^\star_{16,4},
\]
so the final-weight constraint is active and the maximizing value is
$\zeta=\delta$.

Conditional on a block being active, the optimizing local-support distribution
is approximately
\[
\begin{array}{c|c}
\text{local support} & \text{conditional fraction}\\
\hline
13 & 0.00152\\
14 & 0.02159\\
15 & 0.19035\\
16 & 0.78653
\end{array}
\]
Thus the transition is dominated by dense local shortenings rather than
minimum-weight words.

In bits per output coordinate, the exponent decomposes approximately as
\[
    \cS_{16,4}\approx+0.1319884,
    \qquad
    \Gamma_1\approx-0.0681350,
    \qquad
    \Gamma_2\approx-0.0638533,
\]
which sum to zero at the transition.  This is about $81.3\%$ of the
large-field rate-$1/4$ GV relative distance.

\paragraph{Sparse tail and finite-length margins.}

Here $d=13$, so the one-active-block structural tail has scale
\[
    \boxed{N^{-12}.}
\]
The leading marked term is
\[
    \boxed{
    \mu_1^{\rm mark}
      \sim
    \frac{13!}{16}
    \binom{16}{13}^{\!2}
    \left(\frac{\delta}{1-\delta}\right)^{13}
    N^{-12}.
    }
    \tag{R6}
\]
At output length $N=2^{20}$, the numerical estimates are
\[
\begin{array}{c|c}
\delta & -\log_2\mu_1^{\rm mark}\\
\hline
0.55   & 189.44\\
0.58   & 187.15\\
0.60   & 185.60\\
0.6044 & 185.26
\end{array}
\]
The estimated linear-weight margins at $N=2^{20}$ are
\[
\begin{array}{c|c|c}
\delta
& \max\Phi_{16,4,\delta}/\log 2
& -N\max\Phi_{16,4,\delta}\text{ (bits)}\\
\hline
0.6000  & -0.0168867 & 17707\\
0.6040  & -0.0015990 & 1677\\
0.6042  & -0.0008257 & 866\\
0.6043  & -0.0004387 & 460\\
0.60435 & -0.0002451 & 257\\
0.6044  & -0.0000515 & 54
\end{array}
\]
For $N=2^{20}$, evaluating the first moment at
\[
    \boxed{
    C_0=\operatorname{RS}[16,4,13],
    \qquad
    \delta=0.6043
    }
    \tag{R7}
\]
gives a numerical margin of roughly $40$ bits.  This does not certify that
$d_{\min}>0.6043N$.  At $\delta=0.60$, the estimated margin is much larger.

\subsection{Structural transitions}

Entropy optimization of (G9) gives the following estimated transitions:
\[
\begin{array}{c|cc}
\text{scheme} & \text{trace first moment} & \text{direct structural container}\\
\hline
\operatorname{RS}[16,8,9] & 0.3052795 & 0.3084210\\
\operatorname{RS}[16,4,13] & 0.6044133 & 0.6093755
\end{array}
\tag{G10}
\]
At the trace thresholds, the direct-container exponents remain approximately
$-0.01101$ and $-0.02115$ bits per coordinate, respectively.  Thus trace
clustering creates real slack, but the recovered relative distance is only
about $0.0031$ and $0.0050$.

The rate-$1/2$ structural saddle has normalized statistics
\[
\begin{aligned}
|T|/N&\approx0.308421,& r_1/N&\approx0.100975,\\
|\partial T|/N&\approx0.409396,& r_2/N&\approx0.120308,\\
|U(T)|/N&\approx0.529704,& \kappa(U(T))/N&\approx0.308421.
\end{aligned}
\tag{G11}
\]
At rate $1/4$ the corresponding values are
\[
\begin{aligned}
|T|/N&\approx0.609375,& r_1/N&\approx0.102787,\\
|\partial T|/N&\approx0.712162,& r_2/N&\approx0.069877,\\
|U(T)|/N&\approx0.782039,& \kappa(U(T))/N&\approx0.609375.
\end{aligned}
\tag{G12}
\]
Moderate-length integer-program experiments on the exact min-cut formulation
produce block-saturated witnesses consistent with these profiles.  These
experiments and all numbers in (G10)--(G12) are diagnostic rather than theorem
statements.

A positive exponent in (G9) is not a typical-code upper bound without a
concentration argument for structural containers.  Still, the
coefficient-independent transition lies only $0.0031$ or $0.0050$ beyond the
trace threshold.  The numerical gap to $q$-ary GV is much larger.

\subsection{Comparison across alphabets}
\label{sec:comparison}

Table~\ref{tab:comparison} places the first-moment and structural transitions
on one scale.  Distances use the natural symbol Hamming metric for each
alphabet.

\begin{table}[ht]
\centering
\scriptsize
\resizebox{\textwidth}{!}{%
\begin{tabular}{@{}llp{3.1cm}cccc@{}}
\toprule
Alphabet & Rate & Local code & Trace/enumerator $\delta^\star$ & Rank-container $\delta_{\rm str}$ & GV & Sparse scale \\
\midrule
$\F_2$ & $1/2$ & Golay $[24,12,8]$ & $0.105119$ & -- & $0.110028$ & $N^{-3}$ \\
$\F_2$ & $1/2$ & $\operatorname{RM}(2,5)$ $[32,16,8]$ & $0.106165$ & -- & $0.110028$ & $N^{-3}$ \\
large $\F_q$ & $1/2$ & RS $[16,8,9]$ & $0.305280$ & $0.308421$ & $0.492189$ & $N^{-8}$ \\
\addlinespace
$\F_2$ & $1/4$ & BCH subcode $[32,8,12]$ & $0.213293$ & -- & $0.214502$ & $N^{-5}$ \\
large $\F_q$ & $1/4$ & RS $[16,4,13]$ & $0.604413$ & $0.609376$ & $0.743584$ & $N^{-12}$ \\
\bottomrule
\end{tabular}}
\caption{Numerical transitions.  Binary rows use the exact weight-enumerator
first moment.  Large-field trace rows use the support/shortening first moment,
while rank-container rows count the coefficient-independent obstruction.  The
$q$-ary GV values use $q=2^{128}$.}
\label{tab:comparison}
\end{table}

At equal rate, large-field BAA has about $2.8$ times the relative symbol
distance of binary BAA and smaller sparse code-generation tails.  Relative to
the alphabet-specific benchmark, however, the binary schemes are much closer
to GV.  The $\operatorname{RM}(2,5)$ and $[32,8,12]$ schemes reach $96.5\%$ and
$99.4\%$ of binary GV.  The large-field trace transitions reach $62.0\%$ and
$81.3\%$ of their $q$-ary GV values; the structural thresholds improve these
fractions only to $62.7\%$ and $81.9\%$.

The difference comes from scalar freedom.  Over $\F_2$, support parity fixes
the accumulator state.  Over a large field, an $r$-dimensional shortening has
about $q^{r-1}$ projective candidates and can absorb as many as $r-1$
homogeneous cancellation equations without a $1/q$ penalty.  Projectivization
removes duplicate scalar words, but it does not remove block saturation.

\section{Implications and open problems}
\label{sec:implications}

\subsection{Limits of scalar BAA}

Numerical optimization of the trace exponent gives diminishing improvements as
the rate-matched MDS block grows.

\paragraph{Rate $1/2$.}
\[
\begin{array}{c|c}
\text{local block} & \text{estimated trace transition}\\
\hline
{[8,4,5]}    & 0.2847\\
{[12,6,7]}   & 0.2987\\
{[16,8,9]}   & 0.3053\\
{[20,10,11]} & 0.3091\\
{[24,12,13]} & 0.3116\\
{[32,16,17]} & 0.3146\\
{[64,32,33]} & 0.3191
\end{array}
\]
The numerical large-block trend is approximately $0.3234$.

\paragraph{Rate $1/4$.}
\[
\begin{array}{c|c}
\text{local block} & \text{estimated trace transition}\\
\hline
{[8,2,7]}    & 0.5884\\
{[12,3,10]}  & 0.5994\\
{[16,4,13]}  & 0.6044\\
{[20,5,16]}  & 0.6073\\
{[24,6,19]}  & 0.6092\\
{[32,8,25]}  & 0.6115\\
{[64,16,49]} & 0.6147
\end{array}
\]
The numerical large-block trend is approximately $0.618$.

Larger blocks reduce block-activation entropy but do not eliminate two-hop
block saturation.  An MDS code already minimizes every shortening dimension
among codes with the same $(m,k)$.  The data therefore give no indication that
a different local linear code, or a much larger rate-matched MDS block, closes
the gap to GV within scalar BAA.

\paragraph{Architectural consequences.}

The two numerical bounds are summarized by
\[
\boxed{
\begin{array}{c|ccc}
\text{rate} & \text{trace threshold} & \text{structural region} & q\text{-ary GV}\\
\hline
1/2 & 0.3053 & \approx0.3084 & 0.4922\\
1/4 & 0.6044 & \approx0.6094 & 0.7436
\end{array}}
\tag{C1}
\]
The structural value is not a rigorous converse because concentration is still
open.  Its proximity to the trace threshold nevertheless shows that trace
multiplicity accounts for only a small part of the gap to GV.

Binary chosen-block BAA behaves differently.  Support parity fixes the
accumulator state, so there is no projective cancellation budget or large-field
block-rank family.  The large-field encoder still has greater absolute symbol
distance, $O(N)$ work, and sparse tails of order $N^{-8}$ or $N^{-12}$.  The
rank calculation, however, indicates that approaching full $q$-ary GV requires
a different recursive stage.

The obstruction suggests four modifications:
\begin{enumerate}
    \item adding a third independent monomial-interleaver/accumulator layer;
          multiple-accumulator coding theorems give precedent for closing GV
          gaps in binary RMA ensembles~\cite{kliewer08};
    \item encoding at lower rate and puncturing to the target rate, again
          following classical RMA methodology~\cite{kliewer08};
    \item replacing the scalar accumulator by a constant-state vector recursion
          in which a reset consumes more than one independent field equation;
    \item designing a recursive layer whose backward graph expands across local
          blocks rather than allowing two-hop neighborhoods to saturate them.
\end{enumerate}
Each modification permits $O(N)$ encoding when the state dimension is constant.
Unlike a sharper union bound, each also changes the structural cut problem.

\subsection{Open problems}

At $q=2^{128}$ and $N=2^{20}$, the two candidate finite-length bounds are
\[
    d_{\min}>0.30N
    \quad\text{for }\operatorname{RS}[16,8,9],
\]
and
\[
    d_{\min}>0.6043N
    \quad\text{for }\operatorname{RS}[16,4,13],
\]
Establishing either inequality requires interval certificates for the trace
exponents, uniform bounds for sparse and mesoscopic weights, and exact
finite-$N$ prefactors.

On the structural side, the principal open problem is concentration of the
container count $Z_t$ around the transition in (G10), including the possibility
of cheaper cuts that use intermediate-layer vertices rather than only outer
block capacity.  A matching upper and lower analysis would turn the numerical
region into an ensemble-distance characterization.

The broader construction question is whether a slightly more complex
linear-time recursion can retain BAA's implementation advantages while moving
the generic-rank min-cut toward the $q$-ary GV/Singleton frontier.

\section{Formal verification and proof audit}
\label{sec:lean}

The accompanying Lean 4 development formalizes the finite algebraic and
combinatorial results.  Its purpose is to audit the proof at the points where
informal notation can hide a boundary case: sample spaces, index ranges,
endpoint conventions, and shortening hypotheses.  The artifact is contained
in the \texttt{lean/} directory and is pinned to Lean and mathlib
version~\texttt{v4.32.1}.

\subsection{Representation and proof architecture}

A length-$n$ word over $F$ is represented as a function
$\operatorname{Fin}(n)\to F$, and a linear code is a submodule of that vector
space.  Hamming weight uses mathlib's \texttt{hammingNorm}.  The accumulator,
difference, permutation, diagonal, and monomial transformations are linear
equivalences.  The repeated outer code uses an explicit equivalence between
$\operatorname{Fin}(B m)$ and
$\operatorname{Fin}(B)\times\operatorname{Fin}(m)$, which removes hidden
divisibility and coordinate-order assumptions.

Ensemble probabilities are represented as ratios of finite cardinalities in
$\mathbb Q$.  Both first-moment arguments are therefore exact finite counts;
no measure-theoretic layer is needed.  Generic rank is defined over the fraction
field of a multivariate polynomial ring.  The two diagonal layers use disjoint
variable families, so different Cauchy--Binet selections yield different
monomials.

The concrete binary codes also have a computation-oriented representation based
on fixed-width bit vectors.  Conversion lemmas connect these evaluators to the
abstract submodules.  Consequently, a computed spectrum certifies the code in
the corresponding mathematical theorem.

\subsection{Verified statements}

The development contains complete proofs of the following finite claims.

\begin{itemize}
    \item The BAA transformations are invertible and preserve dimension.  The
          block-direct-sum enumerator and the binary identities (B1)--(B5) are
          exact finite theorems.
    \item The extended Golay and $\operatorname{RM}(2,5)$ constructions have
          the dimensions, minimum distances, and spectra stated above.  The
          eight-row $[32,8,12]$ code has the stated spectrum and lies in a
          concretely defined extended primitive BCH code.
    \item Reed--Solomon evaluation codes have the claimed dimension and
          minimum distance.  Their MDS shortening dimensions give (5), (L1),
          (L2), (R3), and (R4).
    \item The zero-return placement count (1), projective bound (2), master
          bound (3), and corrected support-excess bound (4) are formal
          consequences of finite counting and linear algebra.
    \item The backward parity-check identities (G1)--(G5), run and block-rank
          counts (G6)--(G8), and endpoint-refined expectation (G9) are proved.
          The generic rank equals the maximum size of a two-layer
          vertex-disjoint path packing with the local MDS capacities.
\end{itemize}

The path-packing result is stronger than the neighborhood upper bound used in
(G5).  Its proof first establishes a rectangular Cauchy--Binet identity.  It
then identifies nonzero minors of the difference matrix with unique
self-or-successor linkages.  Finally, the MDS parity-check theorem converts
selected columns in each local block into the capacity
$\min\{|U\cap B|,m-k\}$.

\subsection{Audit findings and trust boundary}

Formalization required three corrections to the informal statements.

\begin{enumerate}
    \item Formula (B2) requires $h\ge1$; a separate theorem handles the
          impossible output weight $h=0$.
    \item The Singleton step before (4) requires every touched block to have a
          nontrivial local shortening.  Exact supports of outer codewords meet
          this condition, which is the case used later.
    \item The former version of (G9) suppressed two endpoint corrections while
          displaying a finite equality.  Equation~(G9) now uses flags
          $e_1,e_2$ and is exact.  Suppressing the flags remains valid only at
          the entropy-exponent level.
\end{enumerate}

None of the corrections changes the BAA construction, the finite master bound,
or the reported asymptotic exponents.  No contradiction was found within the
formalized finite scope.  The analytic and numerical claims lie outside that
scope.

\paragraph{Trust boundary.}

The artifact contains no \texttt{sorry}, \texttt{admit}, or project-local
axiom.  A dedicated module applies \texttt{\#print axioms} to the theorems
stated in the paper.  Those theorems depend only on the standard Lean/mathlib
foundations.  The audit reports \texttt{propext},
\texttt{Classical.choice}, and \texttt{Quot.sound}.  Native decision procedures
are used only for labeled small-instance regression tests.  No theorem stated
in the paper depends on a generated native-decision axiom.  The fixed-code
spectra and distances have kernel-checkable proofs.

The command \texttt{lake build BAA.AxiomAudit}, run from \texttt{lean/}, checks
the complete development and prints the dependency audit.  Exhaustive
regressions cover the binary accumulator and difference-linkage conventions
through length eight and the nonbinary placement formula through length seven.
These tests supplement the symbolic all-length theorems; they do not replace
them.

The Lean development does not claim the one-accumulator asymptotic limit,
(B6)--(B12), (6)--(7), (L3)--(L7), any numerical transition, the $O(N)$ cost
model, or concentration of structural containers.  These claims require
analytic estimates, certified optimization, a computational cost model, or new
probabilistic arguments.  Separating them from the verified finite core makes
the remaining proof obligations explicit.

\section{Conclusion}

Chosen-block BAA replaces a random local outer map by a reusable, certified
code.  In the binary setting, this yields exact finite first moments and
near-GV numerical transitions with simple linear-time encoding.  Over large
fields, support-based counting avoids the misleading multiplicity of scalar
words, but generic rank reveals a separate obstruction caused by saturated
local blocks after two backward difference expansions.

Lean checks the finite algebraic and combinatorial claims and accounts for the
three corrections described in Section~\ref{sec:lean}.  The remaining problems
are analytic and probabilistic: certifying the numerical optimizations,
controlling sparse and mesoscopic weights uniformly, and proving or disproving
concentration of structural containers.  For large-field codes near GV, the
rank obstruction points toward a different recursion rather than another bound
for the same scalar two-layer construction.

\appendix
\renewcommand{\theHsection}{appendix.\Alph{section}}
\section{Explicit binary [32,8,12] subcode}
\label{app:binary-generator}

The following eight $32$-bit rows generate the rate-$1/4$ binary block used in
Section~\ref{sec:binary}.  Hexadecimal words are interpreted with the
least-significant bit as coordinate zero:
\[
\begin{array}{cccc}
\texttt{FFFFFFFF} & \texttt{B7DC4AE7} & \texttt{AA9D2399} & \texttt{60E73723}\\
\texttt{6844CC86} & \texttt{180D3C68} & \texttt{52610E07} & \texttt{19DE6790}.
\end{array}
\]
The rows are linearly independent.  Exhaustive enumeration of the resulting
$2^8$ words gives minimum distance $12$ and the exact spectrum
\[
    1+48z^{12}+158z^{16}+48z^{20}+z^{32}.
\]
The code is a subcode of the extended primitive BCH code $[32,11,12]$.

\begin{thebibliography}{99}

\bibitem{baa26}
Vladimir Kolesnikov, Stanislav Peceny, Rahul Rachuri, Srinivasan Raghuraman,
Peter Rindal, and Harshal Shah.
\newblock Block-Accumulate Codes: Accelerated Linear Codes for PCGs and ZK.
\newblock In \emph{CRYPTO 2026, Part VIII}, LNCS 16807, pages 397--429, 2026.
Also IACR ePrint 2025/1828.

\bibitem{berrou93}
Claude Berrou, Alain Glavieux, and Punya Thitimajshima.
\newblock Near Shannon limit error-correcting coding and decoding: Turbo codes.
\newblock In \emph{Proceedings of ICC '93}, pages 1064--1070, 1993.

\bibitem{benedetto98}
Sergio Benedetto, Dariush Divsalar, Guido Montorsi, and Fabrizio Pollara.
\newblock Serial concatenation of interleaved codes: Performance analysis,
design, and iterative decoding.
\newblock \emph{IEEE Transactions on Information Theory}, 44(3):909--926, 1998.

\bibitem{divsalar98}
Dariush Divsalar, Hui Jin, and Robert J. McEliece.
\newblock Coding theorems for ``turbo-like'' codes.
\newblock In \emph{36th Allerton Conference on Communication, Control, and
Computing}, pages 201--210, 1998.

\bibitem{kahale98}
Nabil Kahale and Ruediger Urbanke.
\newblock On the minimum distance of parallel and serially concatenated codes.
\newblock In \emph{IEEE International Symposium on Information Theory}, page
31, 1998.

\bibitem{bms09}
Louay Bazzi, Mohammad Mahdian, and Daniel A. Spielman.
\newblock The minimum distance of turbo-like codes.
\newblock \emph{IEEE Transactions on Information Theory}, 55(1):6--15, 2009.

\bibitem{kliewer08}
Joerg Kliewer, Kamil S. Zigangirov, Christian Koller, and Daniel J. Costello, Jr.
\newblock Coding theorems for repeat multiple accumulate codes.
\newblock arXiv:0810.3422, 2008.

\bibitem{yang04}
Kyeongcheol Yang.
\newblock Weighted nonbinary repeat-accumulate codes.
\newblock \emph{IEEE Transactions on Information Theory}, 50(3):527--531,
2004.

\bibitem{rosnes11}
Eirik Rosnes and Alexandre Graell i Amat.
\newblock On the analysis of weighted nonbinary repeat multiple-accumulate
codes.
\newblock arXiv:1101.5966, 2011.

\bibitem{akhiani26}
Pariya Akhiani and Yupeng Zhang.
\newblock Distance of RAA codes over large finite fields (with applications in
zkSNARKs and PCGs).
\newblock IACR ePrint 2026/524, 2026.

\bibitem{khabbazian26}
Majid Khabbazian.
\newblock Sharp minimum-distance lower tails for RAA codes.
\newblock IACR ePrint 2026/1075, 2026.

\bibitem{bcgi18}
Elette Boyle, Geoffroy Couteau, Niv Gilboa, and Yuval Ishai.
\newblock Compressing vector OLE.
\newblock In \emph{ACM CCS 2018}, pages 896--912, 2018.

\bibitem{bcg19pcg}
Elette Boyle, Geoffroy Couteau, Niv Gilboa, Yuval Ishai, Lisa Kohl, and Peter
Scholl.
\newblock Efficient pseudorandom correlation generators: Silent OT extension
and more.
\newblock In \emph{CRYPTO 2019, Part III}, LNCS 11694, pages 489--518, 2019.

\bibitem{bcg20ring}
Elette Boyle, Geoffroy Couteau, Niv Gilboa, Yuval Ishai, Lisa Kohl, and Peter
Scholl.
\newblock Efficient pseudorandom correlation generators from Ring-LPN.
\newblock In \emph{CRYPTO 2020, Part II}, LNCS 12171, pages 387--416, 2020.

\bibitem{silver21}
Geoffroy Couteau, Peter Rindal, and Srinivasan Raghuraman.
\newblock Silver: Silent VOLE and oblivious transfer from hardness of decoding
structured LDPC codes.
\newblock In \emph{CRYPTO 2021, Part III}, LNCS 12827, pages 502--534, 2021.

\bibitem{ea22}
Elette Boyle, Geoffroy Couteau, Niv Gilboa, Yuval Ishai, Lisa Kohl, Nicolas
Resch, and Peter Scholl.
\newblock Correlated pseudorandomness from Expand-Accumulate codes.
\newblock In \emph{CRYPTO 2022, Part II}, LNCS 13508, pages 603--633, 2022.

\bibitem{ec23}
Srinivasan Raghuraman, Peter Rindal, and Titouan Tanguy.
\newblock Expand-Convolute codes for pseudorandom correlation generators from
LPN.
\newblock In \emph{CRYPTO 2023, Part IV}, LNCS 14084, pages 602--632, 2023.

\bibitem{stationary25}
Vladimir Kolesnikov, Stanislav Peceny, Srinivasan Raghuraman, and Peter Rindal.
\newblock Stationary syndrome decoding for improved PCGs.
\newblock In \emph{CRYPTO 2025}, pages 284--317, 2025.  Also IACR ePrint
2025/295.

\bibitem{lxyy25}
Zhe Li, Chaoping Xing, Yizhou Yao, and Chen Yuan.
\newblock Efficient pseudorandom correlation generators for any finite field.
\newblock IACR ePrint 2025/169, 2025.

\bibitem{rr25}
Noga Ron-Zewi and Ron D. Rothblum.
\newblock Proving as fast as computing: Succinct arguments with constant prover
overhead.
\newblock \emph{Journal of the ACM}, 72(2):15:1--15:54, 2025.
DOI 10.1145/3721477.

\bibitem{orion22}
Tiancheng Xie, Yupeng Zhang, and Dawn Song.
\newblock Orion: Zero knowledge proof with linear prover time.
\newblock In \emph{CRYPTO 2022}, pages 299--328, 2022.  Also IACR ePrint
2022/1010.

\bibitem{brakedown23}
Alexander Golovnev, Jonathan Lee, Srinath Setty, Justin Thaler, and Riad S.
Wahby.
\newblock Brakedown: Linear-time and field-agnostic SNARKs for R1CS.
\newblock In \emph{CRYPTO 2023, Part II}, LNCS 14082, pages 193--226, 2023.

\bibitem{fieldEA24}
Alexander R. Block, Zhiyong Fang, Jonathan Katz, Justin Thaler, Hendrik Waldner,
and Yupeng Zhang.
\newblock Field-agnostic SNARKs from Expand-Accumulate codes.
\newblock In \emph{CRYPTO 2024}, LNCS 14929, pages 276--307, 2024.

\bibitem{blaze25}
Martijn Brehm, Binyi Chen, Ben Fisch, Nicolas Resch, Ron D. Rothblum, and Hadas
Zeilberger.
\newblock Blaze: Fast SNARKs from interleaved RAA codes.
\newblock In \emph{EUROCRYPT 2025}, pages 123--152, 2025.  Also IACR ePrint
2024/1609.

\bibitem{bolt26}
Kobi Gurkan, Andrija Novakovic, and Ron D. Rothblum.
\newblock Bolt: Faster SNARKs from sketched codes.
\newblock In \emph{CRYPTO 2026, Part IX}, pages 362--395, 2026.  Also IACR
ePrint 2026/310.

\enlargethispage{3\baselineskip}
\bibitem{pipfri26}
Weihan Li, Zongyang Zhang, Sherman S. M. Chow, Yanpei Guo, Boyuan Gao, Xuyang
Song, Yi Deng, and Jianwei Liu.
\newblock Shred-to-shine metamorphosis of (distributed) polynomial commitments.
\newblock In \emph{USENIX Security 2026}, 2026.

\bibitem{macwilliams}
F. J. MacWilliams and N. J. A. Sloane.
\newblock \emph{The Theory of Error-Correcting Codes}.
\newblock North-Holland, 1977.

\end{thebibliography}

\end{document}
